%%%%%%%% CS6140 FINAL PROJECT REPORT — ICML 2026 FORMAT %%%%%%%%
%
% SETUP: Download icml2026.sty from the official ICML 2026 style files:
%   https://media.icml.cc/Conferences/ICML2026/Styles/icml2026.zip
% Extract and place icml2026.sty (and fancyhdr.sty if needed) in the
% same directory as this file.
%
% This template is adapted from the ICML 2026 example paper for use as
% a CS6140 Machine Learning course project report. The double-blind
% requirement has been removed (we use the [accepted] option).
%
% COMPILATION: Use pdflatex + bibtex, or set Overleaf to use pdfLaTeX.
%
%%%%%%%%%%%%%%%%%%%%%%%%%%%%%%%%%%%%%%%%%%%%%%%%%%%%%%%%%%%%%%%%%%%%%%%%%

\documentclass{article}

% Recommended packages
\usepackage{microtype}
\usepackage{graphicx}
\usepackage{subfig}
\usepackage{booktabs}       % For professional tables
\usepackage{amsmath}
\usepackage{amssymb}

\usepackage{hyperref}

% Use [accepted] to show author names (this is a course project, not blind review)
\usepackage[accepted]{icml2026}

% Optional: additional packages you may find useful
% \usepackage{algorithm}
% \usepackage{algorithmic}
% \usepackage{listings}

%%%%%%%%%%%%%%%%%%%%%%%%%%%%%%%%%%%%%%%%%%%%%%%%%%%%%%%%%%%%%%%%%%%%%%%%%
% FILL IN YOUR DETAILS BELOW
%%%%%%%%%%%%%%%%%%%%%%%%%%%%%%%%%%%%%%%%%%%%%%%%%%%%%%%%%%%%%%%%%%%%%%%%%

\icmltitlerunning{CS6140 Final Project — Your Short Title Here}

\begin{document}

\twocolumn[
\icmltitle{Your Project Title Here}

% List all team members. Add or remove \icmlauthor blocks as needed.
\icmlsetsymbol{equal}{*}

\begin{icmlauthorlist}
\icmlauthor{Team Member 1}{neu}
\icmlauthor{Team Member 2}{neu}
\icmlauthor{Team Member 3}{neu}
\end{icmlauthorlist}

\icmlaffiliation{neu}{Khoury College of Computer Sciences, Northeastern University, Portland, ME}

\icmlcorrespondingauthor{Team Member 1}{lastname.f@northeastern.edu}

\icmlkeywords{Machine Learning, CS6140}

\vskip 0.3in
]

\printAffiliationsAndNotice{}

%%%%%%%%%%%%%%%%%%%%%%%%%%%%%%%%%%%%%%%%%%%%%%%%%%%%%%%%%%%%%%%%%%%%%%%%%
% ABSTRACT
%%%%%%%%%%%%%%%%%%%%%%%%%%%%%%%%%%%%%%%%%%%%%%%%%%%%%%%%%%%%%%%%%%%%%%%%%

\begin{abstract}
% Write a 150–250 word summary of your project. Include:
% - The problem you are solving
% - Your approach (baseline + advanced model)
% - Key results
% - Main takeaway
%
% TIP: Write this last, after all other sections are complete.
\end{abstract}

%%%%%%%%%%%%%%%%%%%%%%%%%%%%%%%%%%%%%%%%%%%%%%%%%%%%%%%%%%%%%%%%%%%%%%%%%
% 1. INTRODUCTION — due with Update 1 (Literature Review)
%%%%%%%%%%%%%%%%%%%%%%%%%%%%%%%%%%%%%%%%%%%%%%%%%%%%%%%%%%%%%%%%%%%%%%%%%

\section{Introduction}
\label{sec:introduction}

% Describe the problem you are trying to solve.
% - Use formal language to define the inputs and outputs.
% - Why does this problem matter? Who benefits from a solution?
% - What makes this problem challenging?
% - Briefly preview your approach and results.

%%%%%%%%%%%%%%%%%%%%%%%%%%%%%%%%%%%%%%%%%%%%%%%%%%%%%%%%%%%%%%%%%%%%%%%%%
% 2. RELATED WORK — due with Update 1 (Literature Review)
%%%%%%%%%%%%%%%%%%%%%%%%%%%%%%%%%%%%%%%%%%%%%%%%%%%%%%%%%%%%%%%%%%%%%%%%%

\section{Related Work}
\label{sec:related-work}

% Survey the relevant literature (aim for 8–15 references).
% - What methods have others used for this or similar problems?
% - What are the strengths and limitations of prior approaches?
% - How does your work differ from or build on existing work?
%
% TIP: Organize by theme or method family, not as a list of paper summaries.

%%%%%%%%%%%%%%%%%%%%%%%%%%%%%%%%%%%%%%%%%%%%%%%%%%%%%%%%%%%%%%%%%%%%%%%%%
% 3. DATA — due with Update 2 (Data & Methodology)
%%%%%%%%%%%%%%%%%%%%%%%%%%%%%%%%%%%%%%%%%%%%%%%%%%%%%%%%%%%%%%%%%%%%%%%%%

\section{Data}
\label{sec:data}

% Provide a formal definition of your dataset:
% - Source and how it was collected
% - Size (number of samples, features)
% - Feature descriptions and types
% - Distribution of the target variable
% - Missingness and how you handle it
% - Train/validation/test split strategy
% - Any preprocessing or feature engineering steps

% Include at least one table or figure summarizing the data.

%%%%%%%%%%%%%%%%%%%%%%%%%%%%%%%%%%%%%%%%%%%%%%%%%%%%%%%%%%%%%%%%%%%%%%%%%
% 4. METHODS — due with Update 2 (Data & Methodology)
%%%%%%%%%%%%%%%%%%%%%%%%%%%%%%%%%%%%%%%%%%%%%%%%%%%%%%%%%%%%%%%%%%%%%%%%%

\section{Methods}
\label{sec:methods}

% Describe each model you use, with relevant mathematical formulations.

\subsection{Evaluation Metrics}
\label{sec:metrics}

% Define and motivate the evaluation metrics you will use.
% - Why are these metrics appropriate for your problem?
% - Are there trade-offs between metrics (e.g., precision vs. recall)?

\subsection{Baseline Model}
\label{sec:baseline}

% Describe your simple baseline model.
% REQUIREMENT: Must be trainable with scikit-learn.
% - What is the model and why did you choose it?
% - Mathematical formulation (e.g., objective function)
% - Key hyperparameters

\subsection{Advanced Model}
\label{sec:advanced}

% Describe your more complex model.
% REQUIREMENT: Must use deep learning, ensembling, transfer learning,
% or another advanced technique.
% - Architecture or algorithm description
% - Mathematical formulation where appropriate
% - Key hyperparameters

\subsection{Training Procedure}
\label{sec:training}

% Describe your training and validation procedure:
% - Optimization method (SGD, Adam, etc.)
% - Hyperparameter tuning method (grid search, random search, Optuna, etc.)
% - Cross-validation strategy
% - Hardware used (Colab GPU, etc.)

%%%%%%%%%%%%%%%%%%%%%%%%%%%%%%%%%%%%%%%%%%%%%%%%%%%%%%%%%%%%%%%%%%%%%%%%%
% 5. RESULTS — due with Update 3 (Baseline Results) and Final
%%%%%%%%%%%%%%%%%%%%%%%%%%%%%%%%%%%%%%%%%%%%%%%%%%%%%%%%%%%%%%%%%%%%%%%%%

\section{Results}
\label{sec:results}

% Present your results clearly using tables and/or figures.
% - Compare baseline and advanced model performance
% - Report results on your defined evaluation metrics
% - Include error bars or confidence intervals where possible
% - Highlight key findings

% Example results table (uncomment and modify):
% \begin{table}[h]
% \caption{Model performance comparison on the test set.}
% \label{tab:results}
% \centering
% \begin{tabular}{lcc}
% \toprule
% Model & Accuracy & F1 Score \\
% \midrule
% Baseline (Logistic Regression) & 0.XX & 0.XX \\
% Advanced (Your Model) & 0.XX & 0.XX \\
% \bottomrule
% \end{tabular}
% \end{table}

%%%%%%%%%%%%%%%%%%%%%%%%%%%%%%%%%%%%%%%%%%%%%%%%%%%%%%%%%%%%%%%%%%%%%%%%%
% 6. DISCUSSION — due with Final Report
%%%%%%%%%%%%%%%%%%%%%%%%%%%%%%%%%%%%%%%%%%%%%%%%%%%%%%%%%%%%%%%%%%%%%%%%%

\section{Discussion}
\label{sec:discussion}

% Interpret your results:
% - Why does the advanced model outperform (or not) the baseline?
% - What did you learn from the experiments?
% - Error analysis: where does your model fail and why?

\subsection{Limitations}
\label{sec:limitations}

% Be honest about the limitations of your approach:
% - Data limitations (size, bias, representativeness)
% - Model limitations (assumptions, scalability)
% - Evaluation limitations

\subsection{Ethical Considerations}
\label{sec:ethics}

% Discuss potential ethical implications:
% - Could your model cause harm if deployed?
% - Are there fairness or bias concerns in the data or predictions?
% - Privacy considerations?
% - How could these risks be mitigated?

%%%%%%%%%%%%%%%%%%%%%%%%%%%%%%%%%%%%%%%%%%%%%%%%%%%%%%%%%%%%%%%%%%%%%%%%%
% 7. CONCLUSION — due with Final Report
%%%%%%%%%%%%%%%%%%%%%%%%%%%%%%%%%%%%%%%%%%%%%%%%%%%%%%%%%%%%%%%%%%%%%%%%%

\section{Conclusion}
\label{sec:conclusion}

% Summarize your contributions:
% - What problem did you address?
% - What were the key findings?
% - What would you do differently or next with more time?

%%%%%%%%%%%%%%%%%%%%%%%%%%%%%%%%%%%%%%%%%%%%%%%%%%%%%%%%%%%%%%%%%%%%%%%%%
% REFERENCES
%%%%%%%%%%%%%%%%%%%%%%%%%%%%%%%%%%%%%%%%%%%%%%%%%%%%%%%%%%%%%%%%%%%%%%%%%

\bibliography{references}
\bibliographystyle{icml2026}

%%%%%%%%%%%%%%%%%%%%%%%%%%%%%%%%%%%%%%%%%%%%%%%%%%%%%%%%%%%%%%%%%%%%%%%%%
% APPENDIX (optional — does not count toward page limit)
%%%%%%%%%%%%%%%%%%%%%%%%%%%%%%%%%%%%%%%%%%%%%%%%%%%%%%%%%%%%%%%%%%%%%%%%%

% \appendix
% \section{Additional Figures}
% \section{Hyperparameter Search Details}

\end{document}
